\section{Bertini parameters without properness over the field}
\label{sec:kn-relative-nadel-bertini}

\begin{lemma}[A dense smooth parameter open for a nonproper family]
\label{lem:kn-relative-nadel-bertini-good-open}
\lean{EnvelopeProperty.Vanishing.RelativeNadelBertiniGenericOpen.goodOpen_specification}
\leanok
\uses{lem:kn-nadel-bertini-generic-smooth,thm:kn-relative-arithmetic-dense-open-avoiding-closed}
Let $k$ have characteristic zero, let $S$ be an integral affine $k$-scheme,
and let $f:Y\to S$ be a quasi-compact morphism locally of finite presentation.
Suppose that $Y$ is Noetherian and regular.  The open
\[
 U=S\setminus\overline{f(Y\setminus\operatorname{Sm}(f))}
\]
contains the generic point of $S$, is dense and nonempty, and the whole
restriction $f^{-1}(U)\to U$ is smooth.
\end{lemma}
\begin{proof}
The generic fibre is regular: its local rings are localizations of local
rings of $Y$.  The residue field of the generic point has characteristic
zero and is perfect.  Regularity and finite type over that residue field
therefore make the whole generic fibre smooth.  At every point of that
fibre the original morphism is smooth, since the residue-field morphism
at the generic point and its source base change induce isomorphisms on
the relevant local rings.  Thus the image of the nonsmooth locus avoids
the generic point.  The nonsmooth locus is closed, hence constructible
because $Y$ is Noetherian.  Chevalley's theorem makes its image
constructible.  A constructible subset of an integral affine scheme
contains the generic point of its closure whenever that closure is the
whole scheme.  Consequently the closure of this image still avoids the
generic point.  Its complement $U$ contains the generic point and is
dense and nonempty.  Every point of $f^{-1}(U)$ belongs to the smooth
locus, which proves smoothness of the entire restriction.
\end{proof}

\begin{lemma}[Regular equations at actual good parameters]
\label{lem:kn-relative-nadel-bertini-flat-equation}
\lean{EnvelopeProperty.Vanishing.RelativeNadelBertiniFlatSelected.selectedEquation_regular_of_goodOpen}
\leanok
\uses{lem:kn-relative-nadel-bertini-good-open}
Let $C$ be a $k$-algebra, let $(c_i)$ be a finite tuple in $C$ with one
$c_j=1$, and write $h=\sum_i T_ic_i$.  Suppose that the actual quotient
chart $\operatorname{Spec}(C[T]/(h))$ is an open subscheme of a family
$f:Y\to\mathbb A^r_k$.  If a rational parameter $a$ belongs to the
preceding good open of this family, then $\sum_i a_ic_i$ is a
nonzerodivisor of $C$.
\end{lemma}
\begin{proof}
Choose a principal parameter neighbourhood $D(g)$ of $a$ contained in
the good open.  The localized quotient chart is smooth over $k[T]_g$,
hence flat.  The universal equation $h$ is a nonzerodivisor of $C[T]$:
viewing it as a polynomial in $T_j$, its leading coefficient is one.
This remains true after localizing at $g$.  Flatness of the quotient
makes multiplication by the localized equation remain injective after
base change along evaluation at $a$, because the Tor term of the
quotient vanishes.  The original and localized parameter fibres are
canonically isomorphic since $g(a)$ is a unit.  The polynomial fibre
isomorphism identifies the resulting equation with $\sum_i a_ic_i$ in
the original ring $C$.  This proves the claimed injectivity there.
\end{proof}

\begin{lemma}[Relative projective sections give whole Cartier SNC hyperplanes]
\label{lem:kn-relative-nadel-bertini-projective-snc}
\lean{EnvelopeProperty.Vanishing.RelativeNadelBertiniSnc.exists_relative_projective_snc_selected}
\leanok
\uses{lem:kn-relative-nadel-bertini-good-open,lem:kn-relative-nadel-bertini-flat-equation}
Let $X$ be smooth and of finite type over the characteristic-zero field $k$,
and let $j:X\hookrightarrow\mathbb P^m_R$ be a closed immersion for an
arbitrary ring $R$.  For any SNC family $E$ on $X$, a $k$-linear combination
of the original $R$-projective coordinate sections cuts out an effective
Cartier divisor $A$ smooth over $k$, and the whole family $E\cup\{A\}$
is SNC.
\end{lemma}
\begin{proof}
On the inverse image of the $i$-th projective coordinate chart, the
$i$-th coordinate section is an invertible frame.  These inverse images
are affine because $j$ is a closed immersion, and they cover $X$.
In these frames the universal coefficient equation has the form
$\sum_b T_bc_b$ with $c_i=1$.  The incidence projection to $X$ is
therefore smooth: eliminate $T_i$ to identify each chart with an affine
space over that chart of $X$.  Its whole total space is smooth over $k$.
Every closed SNC stratum is also smooth over $k$, since its local ideal
is generated by a subset of the common regular parameters.  Restricting
the original projective embedding to each stratum gives the same chart
argument there.  The preceding good-open lemma thus supplies a
generic-containing coefficient open for the ambient incidence and every
stratum incidence.  There are finitely many strata.  Their intersection
contains a rational $k$-point with a nonzero first coefficient: a nonzero
polynomial over the infinite field $k$ cannot vanish on every tuple.
Choose this one tuple for the ambient space and all strata.

The entire selected zero scheme is smooth, by base change from the smooth
restriction over the coefficient open.  On each original affine chart,
the preceding regular-equation lemma makes its selected equation a
nonzerodivisor, proving that the entire selected ideal is Cartier.
The whole ideal commutes with restriction to each closed stratum because
pullback commutes with the selected section and its zero ideal.  Hence
each such restriction is Cartier and smooth.  At a point of the selected
divisor, quotient by the parameters of all old components through that
point.  The new equation has nonzero cotangent class in the resulting
regular quotient: otherwise quotienting a nonzero Cartier equation
would not give a regular local ring.  Its class is therefore linearly
independent of the old parameter classes.  Complete these classes to
a basis of the cotangent space and lift to regular parameters.  They
simultaneously express all old component ideals and the new Cartier
ideal as distinct coordinate ideals.  At a point outside the new
divisor, the old SNC parameters suffice.  These are exactly the SNC
conditions for the entire appended family.
\end{proof}

\begin{theorem}[Actual whole relative auxiliary families]
\label{thm:kn-relative-nadel-auxiliary-family}
\lean{EnvelopeProperty.Vanishing.RelativeNadelAuxiliaryFamily.exists_veryAmple_auxiliary_family}
\leanok
\uses{lem:kn-relative-nadel-bertini-projective-snc}
Let $X/k$ be smooth of dimension $n$ and of finite type.  Fix an SNC
family $F$, an arbitrary original map $\pi:X\to\operatorname{Spec}R$,
and finitely many lines $M_i$ very ample over that same $R$.  There are
$n+1$ actual sections of each $M_i$ whose whole zero ideals are smooth
Cartier ideals, whose union with the whole original $F$ is SNC, and whose
intersection in each row is empty.
\end{theorem}
\begin{proof}
Use the actual projective embedding and original $\mathcal O(1)$
comparison in each very-ample datum.  The preceding lemma provides a
section transverse to any current finite SNC family.  Apply it
successively to the finite list of all pairs $(i,j)$ with
$0\leq j\leq n$, adjoining each actual whole zero ideal before choosing
the next section.  Transport the sections through the whole original
line comparisons.  This constructs the entire family and proves its
SNC property by induction.  If all $n+1$ divisors in one row contained
a point, their distinct SNC coordinate ideals would supply $n+1$
independent regular parameters in the local ring at that point.  Its
dimension is at most the smooth dimension $n$, which is impossible.
The row intersection is therefore empty as a whole closed subscheme:
its ideal has no point in its support and consequently is the unit ideal.
\end{proof}

\begin{theorem}[Actual whole relative normalized adapted covers]
\label{thm:kn-relative-nadel-cover-geometry}
\lean{EnvelopeProperty.Vanishing.RelativeNadelCoverGeometry.exists_original_finite_flat_smooth_projective}
\leanok
\uses{thm:kn-relative-nadel-auxiliary-family,thm:kn-nadel-auxiliary-positive-powers,lem:kn-nadel-whole-global-cover-finite-smooth}
Assume also that $X$ is integral and let $H$ be an actual ample line
over the original $R$.  For an SNC family $F$ and positive root degrees
$m_i$, some positive power $H^r$ supplies actual sections of
$(H^r)^{m_i}\otimes\mathcal O(-F_i)$ whose auxiliary divisor rows have
empty intersections and whose whole joined family is SNC.  The entire
normalization of the original finite root-cover product is finite flat
over $X$, smooth over $k$ of the original dimension, and projective over
the same original $R$.
\end{theorem}
\begin{proof}
The original relative ample datum and the finitely many invertible
component ideals give a common positive $r$ for which $H^r$ and every
$(H^r)^{m_i}\otimes\mathcal O(-F_i)$ are relatively very ample.  Apply
the preceding theorem to these actual lines.  Compose each obtained
section with the original component-ideal inclusion.  Its whole branch
ideal is the product of that component ideal and its auxiliary zero
ideal.  Thus the original root sections satisfy the actual SNC and
empty-row hypotheses of the normalized adapted-cover construction.
All $m_i$ are units in $k$, by characteristic zero.  The local tame SNC
root calculation makes the entire original normalized cover finite
flat and smooth of the original dimension.  The original $H$ already
makes $X$ projective over $R$.  Finite morphisms preserve projectivity
under composition, so the entire normalized cover remains projective
over that same $R$.
\end{proof}

\begin{theorem}[Actual rational SNC data produces an ample whole relative cover]
\label{thm:kn-relative-nadel-rational-cover-ample}
\lean{EnvelopeProperty.Vanishing.RelativeNadelCoverRationalProduced.exists_original_ample_rounded_cover}
\leanok
\uses{thm:kn-relative-nadel-cover-geometry,lem:kn-nadel-original-rounded-cover-power,thm:kn-relative-finite-ample-original-base}
Let $J=\prod_i F_i^{b_i}$ be an SNC ideal on the original integral
$X/k$ and let $q\geq0$.  Retain an invertible $M$, an integer $d>0$
divisible by $\operatorname{den}(q)$, and actual ampleness over the
original $R$ of $M^d\otimes J^{dq}$.  The original data internally
produces actual auxiliary sections and their whole finite normalized
adapted cover $\nu:W\to X$, smooth over $k$, such that the original
intrinsic rounded root line on $W$ is ample over the same original $R$.
\end{theorem}
\begin{proof}
Use the actual cleared ample line as the seed in the preceding cover
construction, with root degrees $\operatorname{den}(q)$.  This produces
the entire original cover and every original root ideal.  The native
Cartier root-ideal power comparisons give the whole line isomorphism
between the $d$-th power of the rounded root line and
$\nu^*(M^d\otimes J^{dq})$: the original integral rounding and
fractional root weights add to the original rational exponents exactly.
The same comparisons prove that the entire rounded root line is
invertible.  Finite pullback preserves the actual ample criterion over
the original arbitrary ring.  One may see this without changing its
output base by first restricting to the coefficient-image domain of
the integral $X$, applying the whole finite-pullback theorem there,
and composing each actual projective embedding back through the closed
coefficient-base embedding.  The actual projective twist comparison
preserves the original line throughout.  Thus the displayed $d$-th
power is ample over the original $R$, and its original invertible
positive root is ample by its own positive-power criterion.
\end{proof}
